



\section{Frequentist Regret Analysis of \POFUL{}}
\label{sec:regret}
In this section, we present the frequentist regret analysis of POFUL algorithms. Proofs appear in Appendix~\ref{sec:other_proofs}. We first introduce high-probability concentration events.

\begin{definition}
	Fix $\delta \in(0,1)$ and $\delta^\prime = \delta/2T$. We define $\beta^{\PVT}_{t,\delta^\prime,\lambdareg}  \coloneqq \iota_t \beta^{\RLS}_{t,\delta^\prime,\lambdareg} $, $
\hat{\cA}_t \coloneqq \{ \forall s\leq t: \|\widehat{\theta}_t-\theta^{\star}\|_{V_t} \leq \beta_{t,\delta^\prime, \lambdareg}^{RLS} \},\quad
\tilde{\cA}_t  \coloneqq \{ \forall s\leq t: \|\tilde{\theta}_t-\hat{\theta}_t\|_{V_t} \leq \beta^{\PVT}_{t,\delta^\prime,\lambdareg} \}$, and $ \cA_t  \coloneqq \hat{\cA}_t \cap  \tilde{\cA}_t $.

\end{definition}
\begin{proposition}\label{prop:concentration}
		Under Assumptions~\ref{assump:bounded-parameter}, \ref{assump:bounded-action-sets} and \ref{assump:noise} , we have
$\prob{\cA_T} \geq 1 - \delta.
$
\end{proposition}









\subsection{An Data-Driven Regret Bound for \POFUL{}}

In the following, we condition on the event $\cA_T$  which holds with probability $1-\delta$. The following proposition bounds the instantaneous regret of \POFUL{}.
\begin{proposition}\label{prop:instantaneous}
Suppose ${\theta}^\star \in \cE_{t,\delta^\prime, \lambdareg}^{RLS} $ and $\tilde{\theta}_t \in \cE_{t,\delta^\prime, \lambdareg}^{\PVT}$, it holds that 
	\#\label{eq:upperbound}
\dotp{x_t^\star}{ \theta^{\star}}-\dotp{\tilde{x}_t}{\theta^{\star}}\leq (1+\iota_t-\tau_t)\|x^{\star}_t\|_{V^{-1}_t}\beta_{t,\delta^\prime,\lambdareg}^{RLS} +   (1+\iota_t+\tau_t)\|\tilde{x}_t\|_{V^{-1}_t}\beta_{t,\delta^\prime,\lambdareg}^{RLS}.
\#
\end{proposition}


Note that this upper bound is different from what's used in the optimism-based methods \citep{abbasi2011improved,agrawal2013thompson,abeille2017linear}, we reproduce their upper bound and discuss the relationship of our method and theirs in Appendix~\ref{sec:regret_decomp}.



On the right-hand side of \eqref{eq:upperbound}, since the oracle optimal action sequence  $\{x^{\star}_t\}_{t\in[T]}$ is unknown to the algorithm and is different from the action sequence  $\{\tilde{x}_t\}_{t\in[T]}$ played by \POFUL{}, one cannot apply Proposition~\ref{prop:potential} to bound the summation $\sum_{t=1}^T \|\tilde{x}_t\|_{V_t^{-1}}^2 $ and get an upperbound of the regret. To address this,  the key point to connect  $\{\tilde{x}_t\}_{t\in[T]}$ and $\{x^{\star}_t\}_{t\in[T]}$ with $V_t^{-1}$-norm. This motivates the following definition.

\begin{definition}
	For each $t\geq 1$, let $\tilde{x}_t$ and $x_t^\star$ respectively denote the action chosen by \POFUL{} and the optimal action.  We define the \emph{uncertainty ratio} at time $t$ as $\alpha_t \coloneqq \|x_t^\star\|_{V_t^{-1}}/\|\tilde{x}_t\|_{V_t^{-1}}    $. We also define the \emph{(instantaneous) regret proxy} at time $t$ as $\mu_t  \coloneqq  \alpha_t(1+\iota_t -\tau_t)+1+\iota_t+\tau_t$. 

	

\end{definition}

Note that $ \dotp{x}{\widehat{\theta}_t-\theta^{\star}}\leq\|x\|_{V_t^{-1}} \beta_{t,\delta^\prime, \lambdareg}^{RLS}
$ holds with high probability, we have that $\|x\|_{V_t^{-1}} $ essentially determines the length of the confidence interval of the reward $\dotp{x}{\theta^{\star}}$.  Hence, $\alpha_t$ serves as the ratio of uncertainty degrees of the reward obtained by the optimal action $x^{\star}_t$ and the chosen action $\tilde{x}_t$.


 
The intuition behind the definition for $\mu_t$ is constructing a regret upper bound similar to that of OFUL.  Specifically,  Proposition~\ref{prop:instantaneous} indicates
$ \dotp{x_t^\star}{ \theta^{\star}}-\dotp{\tilde{x}_t}{\theta^{\star}} \leq \mu_t  \|\tilde{x}_t\|_{V_t^{-1}} \beta^{RLS}_{t,\delta^\prime, \lambdareg}
$,
and we can check that the instantaneous regret of OFUL satisfies 
$\dotp{x_t^\star}{ \theta^{\star}}-\dotp{\tilde{x}_t}{\theta^{\star}} \leq 2\|\tilde{x}_t\|_{V_t^{-1}}   \beta^{RLS}_{t,\delta^\prime, \lambdareg}$. In this sense, $\mu_t$ is a proxy of the instantaneous regret incurred by \POFUL{} at time $t$. Moreover, OFUL can be regarded as a \POFUL{} algorithm whose $\mu_t$ is fixed at 2, and we could extend the definition of $\alpha_t$ to OFUL by solving $\mu_t  =  \alpha_t(1+\iota_t -\tau_t)+1+\iota_t+\tau_t$ and set $\alpha_t =  1 $ for all $t\in [T]$ for OFUL (recall that in OFUL, $\iota_t = 0$ and $\tau_t = 1$ for all $t \in [T]$). 

% The following Theorem connects $\{\mu_t\}_{t\in[T]}$ and $\cR(T)$ and we postpone the proof to Appendix~\ref{sec:proof_regret}. It provides an oracle but general frequentist regret upper bound for all \POFUL{} algorithms.

The following Theorem connects $\{\mu_t\}_{t\in[T]}$ and $\cR(T)$. It provides an oracle but general frequentist regret upper bound for all \POFUL{} algorithms.


\begin{theorem}[Oracle frequentist regret bound for \POFUL{}] 
\label{thm:oracle_regret} Fix $\delta \in(0,1)$ and let $\delta^\prime = \delta/2T$. Under Assumptions~\ref{assump:bounded-parameter}, \ref{assump:bounded-action-sets} and \ref{assump:noise}, with probability $1-\delta$, 
\POFUL{} achieves a regret of
\#\label{eq:oracle_regret_wo_inflation}
\cR(T) \leq     \sqrt{2 d \left(\sum_{t=1}^T\mu_t^2\right)\log \left(1+\frac{T}{\lambdareg}\right) 
 }\beta_{T,\delta^\prime,\lambdareg}^{RLS}\,.
\#
\end{theorem}

\begin{remark} 
We call Theorem~\ref{thm:oracle_regret} an \emph{oracle} regret bound as $\{\mu_t\}_{t\in[T]}$ for general \POFUL{} depends on the unknown system parameter $\param$. In general, they cannot be calculated by the decision-maker. Nevertheless, note that $\iota_t$ and $\tau_t$ are chosen by the decision-maker, when we have computable upper bounds $\{\hat{\alpha}_t\}_{t\in[T]}$ for $\{\alpha_t\}_{t\in[T]}$, using $\mu_t^2\leq 2\alpha_t^2(1+\iota_t -\tau_t)^2+2(1+\iota_t+\tau_t)^2$, we could calculate upper bounds for $\{\mu_t\}_{t\in[T]}$ as well. Consequently, Theorem~\ref{thm:oracle_regret} instantly turns into a data-driven regret bound for \POFUL{} and could be utilized later for course correction, which will be the aim of the next section. When we additionally know that $1+\iota_t - \tau_t$ is non-negative, we would use the equality $\mu_t = \alpha_t(1+\iota_t -\tau_t)+1+\iota_t+\tau_t$ directly for the bound.
\end{remark}



\begin{remark}
In the Discussion section of \citet{abeille2017linear}, the authors introduce a concept similar to the reciprocal of our $\alpha_t$. They suggest that the necessity of proving LinTS samples are optimistic could be bypassed if for some $\alpha > 0$ LinTS samples $\tilde{\theta}_t$ such that $\|x^\star(\tparam_t)\|_{V_t^{-1}} \geq \alpha \|x^\star(\param_t)\|_{V_t^{-1}}$ with constant probability, where $x^\star(\tparam_t)$ and $x^\star(\param_t)$ represent the optimal actions corresponding to $\tparam_t$ and $\param_t$, respectively. They pose this as an open question regarding the possibility of relaxing the requirement of inflating the posterior. In the following section, we provide a positive answer to this question by studying the reciprocal of their $\alpha$ using geometric arguments. This investigation offers an explanation for the empirical success of LinTS without the need for posterior inflation.
\end{remark}








 



